\subsection{Data}

Our analysis relies on the national panel of school capital elections and outcomes constructed by \citet{biasi_what_2025}. They assemble a unified database of school-bond referenda from administrative records and public information requests, covering local bond elections in 41 U.S. states between 1995 and 2017\footnote{After sample restrictions, their data comprises of 29 states.}. Further, they link these elections to annual district-level finances from the NCES Annual Survey of School System Finances and the Census of Governments, student achievement data from the Stanford Education Data Archive (SEDA), and housing price indices from \citet{contat_larson_2024}. This rich dataset allows them to study the fiscal, academic, and housing market responses to school capital investments authorized through local bond elections. Our analysis uses the cleaned master files in their replication package, and we impose all of their sample restrictions as well as some additional restrictions due to limited sample size for certain states\footnote{Our final analysis consists 24 out of 29 states analysed by \cite{biasi_what_2025}}.

% To study fiscal and demographic responses, \citet{biasi_what_2025} link these elections to annual district-level finances from the NCES Annual Survey of School System Finances and the Census of Governments. These sources provide per-pupil capital expenditures, current instructional and non-instructional expenditures, and revenues by source, beginning in 1995 and expressed in 2020 dollars. District characteristics—enrollment, racial and ethnic composition, and the share of low-income students—come from the NCES Common Core of Data. The resulting panel tracks both fiscal inputs and student composition for virtually all districts that hold at least one capital referendum.

% Student achievement is measured using the extended version of the Stanford Education Data Archive (SEDA) developed in \citet{reardon2017test,fahle2021scaling}. State standardized test scores in math and reading/ELA for grades 3–8 are mapped onto a common national scale using NAEP-based linking. \citet{biasi_what_2025} augment SEDA by incorporating earlier years from the National Longitudinal School-Level State Assessment Score Database and additional state releases, aggregating to district-subject-grade-year means and then standardizing following the SEDA procedure. The final test-score panel covers roughly 10{,}000 districts with at least one bond election, with outcomes expressed in district-level standard deviations.

% To capture housing market responses, \citet{biasi_what_2025} use the house price index constructed by \citet{contat_larson_2024}, a repeat-sales index based on Fannie Mae, Freddie Mac, FHA, and CoreLogic transaction data. The index is available annually for a balanced panel of census tracts from 1989 to 2021 and is normalized to 100 in 1989. Census tract centroids are mapped to 2010 school district boundaries using NCES EDGE shapefiles, and tract-level indices are averaged to obtain a district-level HPI. This yields a balanced sample of 4{,}679 districts with both bond-election and house price information for 1990–2017. In this paper we use the exact district-year panel supplied in the \citet{biasi_what_2025} replication package and do not alter their sample construction, except where noted in our state-level extension.

\subsection{Empirical Framework}

We follow \citet{biasi_what_2025} and exploit quasi-random variation from close bond elections using a dynamic regression discontinuity design. Let $m_{jc}$ denote the vote margin in district $j$ and election cohort $c$, defined as the difference between the fraction of “yes” votes and the required approval threshold in that state. A bond is authorized if and only if $m_{jc} \ge 0$. Under standard RD assumptions, districts where a bond barely passes ($m_{jc} \gtrsim 0$) provide a valid counterfactual for districts where the bond barely fails ($m_{jc} \lesssim 0$), after flexibly controlling for the running variable. The outcomes $Y_{jct}$ are capital expenditures per pupil, standardized test scores, and the house price index, measured at the district-year level.

Because districts can propose and vote on multiple bonds over time, \citet{biasi_what_2025} do not estimate a static, single-election RD. Instead, they develop a \emph{stacked dynamic RD} (DRD) that treats each bond election as a separate “cohort” and follows districts in an event-time window around that election. For each cohort $c$, they keep up to five pre-election years and ten or more post-election years for districts that place a measure on the ballot in year $c$ and either barely pass or barely fail. Finally, they ``stack'' all cohorts into one large panel where an observation is indexed by $(j,c,t)$—district $j$, cohort $c$, and calendar year $t$—and define event time $\tau = t - c$.

In this stacked panel, the main specification relates outcomes to a set of event-time indicators interacted with bond authorization, while controlling for flexible functions of the vote margin and for each district’s bond history. An OLS version of the estimating equation, as presented in \citet{biasi_what_2025}, is

\[
Y_{jct}
= \alpha_{jc} + \gamma_{s(j)t}
+ \sum_{\tau \neq 0} \biggl[\beta_{\tau} \mathbf{1}\{\tau = t-c\} \times \mathbf{1}\{m_{jc} \ge 0\}
+ \phi_{\tau} M_{jct}^{(\tau)}
+ g_{c}(m_{jc})\biggr]
+ u_{jct},
\]

where $\alpha_{jc}$ are district-by-cohort fixed effects, $\gamma_{s(j)t}$ are state-by-year fixed effects, and $g_{c}(m_{jc})$ denotes cohort-specific polynomials in the vote margin. The indicators $M_{jct}^{(\tau)}$ summarize whether district $j$ proposed or passed other bonds $\tau$ years before or after the focal election, so that treated and control districts are compared conditional on their bond histories. The coefficients $\beta_{\tau}$ trace out dynamic treatment effects from five years before to ten or twelve years after authorization with pre-treatment coefficients serve as event-study checks for RD validity. Following \citet{biasi_what_2025}, all regressions are estimated by OLS, weighting by baseline enrollment or the number of test takers for test-score outcomes, and standard errors are clustered at the district level.

In our replication, we implement this stacked DRD specification in Stata using the original Stata code provided by \citet{biasi_what_2025} and port it to R to verify that we can reproduce their tables and figures. For our state-level extension, we use the same estimator and the event-time window, and the same controls for bond histories and vote-margin polynomials, but allow the dynamic treatment effects to vary across states by estimating the model separately for each state.
